\section{Cel i zakres badań}\label{cel-i-zakres-badaux144}

Celem projektu jest ocena wpływu degradacji obrazu na skuteczność
klasyfikacji pojedynczych znaków alfanumerycznych oraz zbadanie, czy
zastosowanie takich degradacji w procesie uczenia zwiększa odporność
niewielkiej sieci konwolucyjnej (CNN). Rozpatrywany jest problem
klasyfikacji wyodrębnionych obrazów znaków, a nie pełny proces
wykrywania i rozpoznawania tekstu w obrazie.

Badanie oparto na porównaniu dwóch modeli o identycznej architekturze:
modelu A uczonego na obrazach bez dodatkowych degradacji oraz modelu B
uczonego z wykorzystaniem syntetycznych zaburzeń. Dla obu modeli
zastosowano wspólny zbiór testowy z tymi samymi poziomami degradacji.

\section{Zbiór danych i przygotowanie
  obrazów}\label{zbiuxf3r-danych-i-przygotowanie-obrazuxf3w}

W eksperymentach wykorzystano EMNIST ByClass, zestaw danych składający
się z 62 klas obejmujących cyfry 0--9, wielkie litery A--Z i małe litery
a--z. Obrazy o rozdzielczości 28 × 28 pikseli skorygowano pod względem
orientacji, po czym uzupełniono dwupikselowym czarnym obramowaniem do 32
× 32 pikseli bez przeskalowywania zawartości. Dane znormalizowano według
średniej 0,1331 i odchylenia standardowego 0,2870.

Z oficjalnego zbioru treningowego (697 933 obrazów) wydzielono przy
stałym ziarnie 123 podzbiór treningowy (628 139) oraz walidacyjny (69
794). Indeksy tego podziału zapisano w plikach projektu. Oficjalny zbiór
testowy liczący 116 323 próbki pozostawiono odrębny od treningu i
walidacji.

\begin{figure}[t]
  \centering
  \includegraphics[width=\linewidth,keepaspectratio]{media/media/image1.png}
  \caption{Przykładowe obrazy i etykiety zbioru EMNIST ByClass po
    korekcie orientacji. Źródło: wyniki własnego przetwarzania danych.}
\end{figure}

\section{Architektura sieci i założenia
  eksperymentu}\label{architektura-sieci-i-zaux142oux17cenia-eksperymentu}

\subsection{Zastosowana architektura
  CNN}\label{zastosowana-architektura-cnn}

Zastosowano własny, kompaktowy klasyfikator z trzema kolejnymi blokami
konwolucyjnymi oraz dwuwarstwową głowicą klasyfikacyjną. Każdy blok
konwolucyjny wykorzystuje filtr 3 × 3 z dopełnieniem zachowującym
rozdzielczość, normalizację wsadową (BatchNorm), funkcję ReLU i operację
MaxPool 2 × 2. Liczba kanałów wzrasta od 32 do 128, natomiast rozmiar
map cech stopniowo maleje od 32 × 32 do 4 × 4.

\begin{longtable}{p{0.24\linewidth}p{0.43\linewidth}p{0.23\linewidth}}
  \toprule
  \textbf{Etap} & \textbf{Operacje}                                   & \textbf{Rozmiar wyjścia} \\
  \midrule
  Wejście       & Obraz w skali szarości                              & 1 × 32 × 32              \\
  Warstwa 1     & Conv 1→32 (3 × 3) + BN + ReLU + MaxPool             & 32 × 16 × 16             \\
  Warstwa 2     & Conv 32→64 (3 × 3) + BN + ReLU + MaxPool            & 64 × 8 × 8               \\
  Warstwa 3     & Conv 64→128 (3 × 3) + BN + ReLU + MaxPool           & 128 × 4 × 4              \\
  Klasyfikator  & Flatten 2048 → Linear 256 + BN + ReLU + Dropout 0,5 & 256                      \\
  Wyjście       & Linear 256 → 62                                     & 62                       \\
  \bottomrule
\end{longtable}

Warstwy konwolucyjne wyznaczają cechy lokalne, a zwiększanie liczby
kanałów umożliwia opis coraz bardziej złożonych kombinacji tych cech.
MaxPool zmniejsza rozdzielczość reprezentacji, ograniczając koszt
obliczeń. Głowica liniowa przekształca uzyskane cechy w 62 wyniki
odpowiadające klasom EMNIST ByClass. Normalizacja wsadowa BatchNorm i
dropout wspierają stabilność uczenia oraz ograniczanie przeuczenia; nie
gwarantują jednak odporności na określone deformacje.

Sieć zawiera 634 110 parametrów uczonych. Same wagi zapisane w formacie
FP32 zajmują około 2,54 MB. Rozmiar sieci stanowi cechę projektu,
natomiast ocena realnej przydatności np. potencjału wybranej
architektury do działania w czasie rzeczywistym na telefonach
komórkowych, wymagałaby osobnych pomiarów na docelowym sprzęcie.

\subsection{Strategie uczenia i degradacje
  obrazu}\label{strategie-uczenia-i-degradacje-obrazu}

Model A trenowano na obrazach bez dodatkowych degradacji. W przypadku
modelu B w trakcie uczenia losowo stosowano trzy typy transformacji;
każda mogła zostać wybrana z prawdopodobieństwem 50\%: rozmycie Gaussa
($\sigma \leq 2{,}0$), szum gaussowski (odchylenie standardowe $\leq 0{,}15$
dla intensywności w zakresie 0--1) oraz zaburzenie perspektywy (maksymalne
przesunięcie narożnika $\leq 0{,}25$ krótszego wymiaru obrazu). Wybranie
jednego z zaburzeń dla danego obrazu nie wyklucza dobrania innego
rodzaju. W efekcie, wszystkie kombinacje 3 możliwych degradacji oraz
braku żadnej znalazły się w zestawie danych treningowych modelu B.

\begin{longtable}{p{0.62\linewidth}p{0.28\linewidth}}
  \toprule
  \textbf{Liczba zaplikowanych degradacji}       & \textbf{Prawdopodobieństwo}              \\
  \midrule
  Wszystkie trzy -- zamazanie, szum, skrzywienie & \textbf{12.5\%} (\({0.5}^{3}\))          \\
  Którekolwiek dwie -- np. zamazanie, szum       & \textbf{37.5\%} (\(3 \times {0.5}^{3}\)) \\
  Tylko jedna -- np. zamazanie                   & \textbf{37.5\%} (\(3 \times {0.5}^{3}\)) \\
  Czysty obraz bez degradacji                    & \textbf{12.5\%} (\({0.5}^{3}\))          \\
  \bottomrule
\end{longtable}

Wybrano trzy różne mechanizmy pogorszenia jakości danych wejściowych:
utratę drobnych szczegółów, zakłócenia wartości pikseli oraz zmianę
geometrii znaku. Na potrzeby wizualizacji przygotowano osobne serie
testów ze stopniowo zwiększaną intensywnością każdego zaburzenia.

\begin{figure}[t]
  \centering
  \includegraphics[width=\linewidth,keepaspectratio]{media/media/image2.png}
  \caption{Przykładowe poziomy rozmycia, szumu i deformacji
    perspektywicznej tego samego znaku. Źródło: własny generator
    degradacji.}
\end{figure}

\section{Przebieg uczenia i
  diagnostyka}\label{przebieg-uczenia-i-diagnostyka}

Oba modele trenowano od początku przez 25 epok, stosując optymalizator
AdamW (learning rate 0,001; weight decay 0,0001), wielkość partii 512 i
funkcję straty entropii krzyżowej. Jako wynik każdego treningu zachowano
checkpoint o najniższej stracie walidacyjnej.

\subsection{Model A --- trening bez
  degradacji}\label{model-a-trening-bez-degradacji}

W modelu A strata treningowa maleje w kolejnych epokach. Strata
walidacyjna osiąga minimum 0,3328 w epoce 11, po czym rośnie do 0,3549 w
epoce 25, mimo dalszej poprawy na zbiorze treningowym. Przebieg ten
wskazuje na pogłębiającą się rozbieżność między dopasowaniem do danych
treningowych a zachowaniem na walidacji. Ostatecznie zapisany model
został pobrany z checkpointu po epoce 11, po której zaobserwowano
zjawisko przeuczenia.

\begin{figure}[t]
  \centering
  \includegraphics[width=\linewidth,keepaspectratio]{media/media/image3.png}
  \caption{Model A: przebieg straty i dokładności w trakcie
    uczenia oraz walidacji.}
\end{figure}

\begin{figure}[t]
  \centering
  \includegraphics[width=\linewidth,keepaspectratio]{media/media/image4.png}
  \caption{Model A: znormalizowana macierz pomyłek na danych
    testowych. Źródło: wyniki własne.}
\end{figure}

\subsection{Model B --- trening z syntetycznymi
  degradacjami}\label{model-b-trening-z-syntetycznymi-degradacjami}

Dla modelu B zadanie treningowe jest trudniejsze, ponieważ część obrazów
w każdej epoce podlega losowym zmianom. W badanym przebiegu strata
walidacyjna maleje do 0,3326 w epoce 25; nie obserwuje się późnego
wzrostu analogicznego do modelu A. W przeciwieństwie do modelu bazowego,
podczas trenowania modelu B nie zaobserwowano zjawiska przeuczenia,
lokalne minimum straty walidacyjnej zostało osiągnięte podczas ostatniej
epoki treningowej.

\begin{figure}[t]
  \centering
  \includegraphics[width=\linewidth,keepaspectratio]{media/media/image5.png}
  \caption{Model B: przebieg straty i dokładności w trakcie
    uczenia oraz walidacji.}
\end{figure}

\begin{figure}[t]
  \centering
  \includegraphics[width=\linewidth,keepaspectratio]{media/media/image6.png}
  \caption{Model B: znormalizowana macierz pomyłek na danych
    testowych. Źródło: wyniki własne.}
\end{figure}

\section{Wyniki testów
  porównawczych}\label{wyniki-testuxf3w-poruxf3wnawczych}

\subsection{Klasyfikacja obrazów bez
  degradacji}\label{klasyfikacja-obrazuxf3w-bez-degradacji}

Na niezmodyfikowanym zbiorze testowym oba modele osiągają zbliżoną
dokładność top-1: 87,62\% dla modelu A i 87,46\% dla modelu B. Różnice
między metrykami uśrednianymi po klasach nie są jednak jednakowe.

\begin{longtable}{p{0.26\linewidth}p{0.22\linewidth}p{0.22\linewidth}p{0.22\linewidth}}
  \toprule
  \textbf{Metryka}                            & \textbf{Model A} & \textbf{Model B} & \textbf{Różnica B-A} \\
  \midrule
  Top-1 accuracy                              & 87,62\%          & 87,46\%          & -0,16 p.p.           \\
  Dokładność bez rozróżniania wielkości liter & 90,96\%          & 90,88\%          & -0,08 p.p.           \\
  Macro-precision                             & 80,59\%          & 82,07\%          & +1,48 p.p.           \\
  Macro-recall                                & 75,30\%          & 74,56\%          & -0,74 p.p.           \\
  Macro-F1                                    & 75,59\%          & 74,58\%          & -1,01 p.p.           \\
  \bottomrule
\end{longtable}

Wprowadzenie degradacji podczas treningu nie spowodowało dużego
pogorszenia dokładności top-1 na obrazach bez degradacji. Jednocześnie
niższy macro-F1 modelu B wskazuje, że trenowanie na obrazach z
degradacjami stanowi kompromis, zamiast być bezwzględnie lepszym
rozwiązaniem.

\subsection{Zmiany dokładności wraz z poziomem
  degradacji}\label{zmiany-dokux142adnoux15bci-wraz-z-poziomem-degradacji}

Oba modele oceniano na tych samych 116 323 próbkach zbioru testowego,
osobno dla kolejnych intensywności degradacji. Poniższe wykresy
przedstawiają pełne przebiegi; zestawienie wybranych punktów ułatwia
porównanie skali różnic.

\begin{longtable}{p{0.33\linewidth}p{0.22\linewidth}p{0.22\linewidth}p{0.20\linewidth}}
  \toprule
  \textbf{Warunek}        & \textbf{Model A} & \textbf{Model B} & \textbf{Różnica B-A} \\
  \midrule
  Rozmycie $\sigma$ = 1,0 & 71,69\%          & 87,25\%          & +15,56 p.p.          \\
  Rozmycie $\sigma$ = 2,0 & 12,24\%          & 86,20\%          & +73,96 p.p.          \\
  Rozmycie $\sigma$ = 3,0 & 4,28\%           & 64,44\%          & +60,16 p.p.          \\
  Szum std = 0,15         & 74,99\%          & 87,08\%          & +12,09 p.p.          \\
  Szum std = 0,20         & 37,65\%          & 86,59\%          & +48,94 p.p.          \\
  Perspektywa = 0,15      & 80,24\%          & 85,83\%          & +5,59 p.p.           \\
  Perspektywa = 0,25      & 64,36\%          & 82,83\%          & +18,47 p.p.          \\
  \bottomrule
\end{longtable}

\subsubsection{Rozmycie gaussowskie}\label{rozmycie-gaussowskie}

Wraz ze wzrostem $\sigma$ dokładność modelu A maleje gwałtownie: od 87,62\% na
czystych obrazach do 12,24\% przy $\sigma$ = 2,0. Model B zachowuje w tym
punkcie 86,20\% dokładności. Przy $\sigma$ = 3,0 wyniki wynoszą odpowiednio
4,28\% i 64,44\%. Poziomy $\sigma$ = 2,5 i 3,0 przekraczają maksymalną wartość
rozmycia stosowaną podczas treningu modelu B.

\begin{figure}[t]
  \centering
  \includegraphics[width=\linewidth,keepaspectratio]{media/media/image7.png}
  \caption{Dokładność top-1 obu modeli w zależności od odchylenia
    standardowego rozmycia gaussowskiego.}
\end{figure}

\subsubsection{Szum gaussowski}\label{szum-gaussowski}

Model B zachowuje zbliżoną dokładność niezależnie od badanej
intensywności szumu. Przy odchyleniu standardowym 0,20 osiąga 86,59\%,
podczas gdy dokładność modelu A spada do 37,65\%. Różnica wynosi 48,94
punktu procentowego.

\begin{figure}[t]
  \centering
  \includegraphics[width=\linewidth,keepaspectratio]{media/media/image8.png}
  \caption{Dokładność top-1 obu modeli w zależności od poziomu
    addytywnego szumu gaussowskiego.}
\end{figure}

\subsubsection{Zniekształcenia
  perspektywiczne}\label{znieksztaux142cenia-perspektywiczne}

Wraz ze wzrostem stopnia skrzywienia dokładność obu modeli stopniowo się
obniża, jednak spadek jest większy dla modelu A. Przy maksymalnym
badanym poziomie 0,25 wartości wynoszą 64,36\% (A) i 82,83\% (B),
różnica wynosi 18,47 punktu procentowego.

\begin{figure}[t]
  \centering
  \includegraphics[width=\linewidth,keepaspectratio]{media/media/image9.png}
  \caption{Dokładność top-1 obu modeli w zależności od
    maksymalnego względnego przesunięcia narożników.}
\end{figure}

\section{Wnioski wstępne i dalszy plan
  badań}\label{wnioski-wstux119pne-i-dalszy-plan-badaux144}

Uzyskane wyniki sugerują, że uwzględnienie syntetycznych degradacji
podczas uczenia zwiększa odporność badanej architektury na rozmycie,
szum i deformacje perspektywiczne, przy niewielkiej zmianie dokładności
na obrazach bez degradacji. Największe przewagi modelu B zaobserwowano
dla silnego rozmycia i szumu. Zakres wnioskowania jest jednak
ograniczony do zastosowanej architektury, zbioru EMNIST ByClass i
przyjętego sposobu generowania degradacji.

\subsection{Ograniczenia}\label{ograniczenia}

Wyniki pochodzą z pojedynczego przebiegu uczenia każdego modelu, dlatego
nie określają jeszcze zmienności między uruchomieniami. Wykorzystano
jedną architekturę i jeden zbiór danych, a degradacje miały charakter
syntetyczny. Badanie nie przesądza więc o skuteczności na zdjęciach
rzeczywistych ani o przenoszeniu obserwacji na inne modele.

\subsection{Planowane kolejne kroki}\label{planowane-kolejne-kroki}

\begin{itemize}
  \item
        Powtórzenie uczenia modeli A i B z różnymi ziarnami aby uniezależnić
        eksperymenty od losowości.
  \item
        Jakościowa analiza pomyłek, w tym par znaków, predykcji obu modeli dla
        tych samych przykładów i wyników według klas.
  \item
        Rozważenie dodatkowych metod degradacji zbliżonych do warunków
        występujących w rzeczywistych warunkach.
  \item
        Ocena możliwości przeniesienia wyników modelu B względem modelu A na
        osobny zbiór rzeczywistych fotografii znaków, przykładowo na zestaw
        danych Chars74k, bądź na samodzielnie stworzony zbiór zdjęć.
\end{itemize}
